\documentclass{practice}

\title{10}
\date{\today}

\usepackage{crysymb}

\begin{document}
\maketitle

\begin{task}{Browsers and revocation}
  Check out \url{https://badssl.com} and the various invalid certificate options on your main browser.
  Test them on other browsers as well for comparison.
  What do you notice?

  Check also \url{https://emudhra.we-have-a-problem.com/#check}.
\end{task}

\begin{task}{Certificate Signing Requests}
  OpenSSL creates and inspects CSRs with the \href{https://docs.openssl.org/master/man1/openssl-req/}{\texttt{req}} tool.
  CSRs are created with the \texttt{-new} option.
  For example, we can generate a P-256 keypair and a CSR with the following command:
  \begin{cmdblock}
openssl req -new -newkey ec -pkeyopt ec_paramgen_curve:P-256 \
  -keyout key.pem \
  -subj "/CN=Jaak-Kristjan Jõeorg/" \
  -out cert.csr \
  -noenc -utf8
  \end{cmdblock}
  \begin{itemize}
    \item \texttt{-newkey} generates a new private key, which is written to the \texttt{-keyout} file.
    If you already have a key, use \texttt{-key} instead.
    \texttt{-noenc} stores the private key unencrypted (not recommended in practice).
    \item \texttt{-utf8} interprets the field values as UTF-8, so that names such as \enquote{Jaak-Kristjan Jõeorg} are encoded correctly.
  \end{itemize}

  If \texttt{-subj} is omitted, OpenSSL prompts for the subject fields interactively.

  \textit{Task.}
  Create a CSR for a P-256 key.
  The subject name must contain exactly the following fields, in this order:
  \begin{itemize}
    \item \texttt{C} (country): \texttt{EE}
    \item \texttt{O} (organisation): \texttt{TalTech}
    \item \texttt{OU} (organisational unit): \texttt{ICS0036}
    \item \texttt{CN} (common name): your full name
    \item \texttt{emailAddress}: your university email address, \texttt{<uni-id>@taltech.ee}
  \end{itemize}

  Submit the CSR to the Discord bot using \texttt{/req}.
\end{task}

\begin{task}{PKCS\#12}
  A certificate on its own is not very useful to its owner: to sign or decrypt, you also need the private key.
  PKCS\#12 (\href{https://datatracker.ietf.org/doc/html/rfc7292}{RFC 7292}) is a container format that bundles a private key, its certificate, and optionally the CA certificates of the chain into a single password-protected file.
  The files usually have the extension \texttt{.p12} or \texttt{.pfx}.

  OpenSSL handles PKCS\#12 files with the \href{https://docs.openssl.org/master/man1/openssl-pkcs12/}{\texttt{pkcs12}} tool.
  To bundle your key and certificate:
  \begin{cmdblock}
openssl pkcs12 -export \
  -inkey key.pem -in cert.pem \
  -name "Friendly name" \
  -out bundle.p12 \
  -iter 600000
  \end{cmdblock}
  where \texttt{-name} optionally sets the \enquote{friendly name} that applications display for the entry.
  CA certificates can be included with \texttt{-certfile chain.pem}.
  By default, OpenSSL only uses 2048 PBKDF2 iterations, so \texttt{-iter} is needed to follow best practices. 

  To view the structure of the file:
  \begin{cmdblock}
openssl pkcs12 -in bundle.p12 -info -noout
  \end{cmdblock}

  \textit{Task.}
  Bundle your key and the certificate into a PKCS\#12 file and inspect it.
\end{task}

\begin{task}{Certificate Revocation Lists}
  OpenSSL inspects CRLs with the \href{https://docs.openssl.org/master/man1/openssl-crl/}{\texttt{crl}} tool.
  \begin{itemize}
    \item Print the CRL in human-readable form:
    \begin{cmdblock}
openssl crl -in issuer.crl -noout -text
    \end{cmdblock}
    \item Verify the CRL signature with the issuer's certificate:
    \begin{cmdblock}
openssl crl -in issuer.crl -noout -CAfile issuer.der
    \end{cmdblock}
  \end{itemize}
\end{task}

\begin{task}{Validating certificate chains}
  OpenSSL validates certificate chains with the \href{https://docs.openssl.org/master/man1/openssl-verify/}{\texttt{verify}} tool.
  It builds a chain from the leaf certificate up to a trust anchor and checks, among others, each signature, the validity periods, and the Basic Constraints and Key Usage extensions of the CA certificates.

  The verifier needs to know which certificates to \emph{trust}, and which certificates may only be \emph{used} to build the chain:
  \begin{itemize}
    \item \texttt{-trusted} gives the trust anchors.
    These are typically root certificates.
    For non-root certificates, \texttt{-partial\_chain} must also be used.
    \item \texttt{-untrusted} gives additional certificates that may be used to build the chain, but are not trusted on their own, e.g. intermediate certificates.
  \end{itemize}
  \texttt{-CAfile} is similar to \texttt{-trusted}, but in addition to the given certificates, OpenSSL also trusts its default trust store.

  \textit{Task.}
  Validate the certificate you received from the bot.
  \begin{enumerate}
    \item Verify your certificate on its own:
    \begin{cmdblock}
openssl verify cert.pem
    \end{cmdblock}

    \item Verify the chain and print it:
    \begin{cmdblock}
openssl verify -show_chain \
  -trusted root.der -untrusted intermediate.der \
  cert.pem
    \end{cmdblock}

    \item The above does not check revocation!
    Download the CRLs from the CRL Distribution Points of your certificate and of the intermediate certificate, and run:
    \begin{cmdblock}
openssl verify -crl_check_all \
  -CRLfile int.crl -CRLfile root.crl \
  -trusted root.der -untrusted intermediate.der \
  cert.pem
    \end{cmdblock}
    Alternatively, \texttt{-crl\_download} downloads the CRLs from the CRL Distribution Points.

    \begin{tcolorbox}[title=NB!]
      You should prefer \texttt{-crl\_check\_all} over \texttt{-crl\_check}.
      The latter option only checks the revocation status of the leaf certificate, so it misses revoked intermediates, even if you provide CRLs that list them (??!).
    \end{tcolorbox}

    \item A valid chain does not yet mean that the certificate fits your use case.
    Try the following options:
    \begin{itemize}
      \item \texttt{-purpose smimesign} and \texttt{-purpose sslserver}
      \item \texttt{-auth\_level 4}
      \item \texttt{-policy 2.23.140.1.5.3.2 -explicit\_policy}
    \end{itemize}
  \end{enumerate}

  Note that by default, OpenSSL tolerates some certificates that do not comply with \href{https://datatracker.ietf.org/doc/html/rfc5280}{RFC 5280}.
  It's good practice to disable these workarounds by using \texttt{-x509\_strict}, but I've omitted it in the example commands for brevity.
\end{task}

\begin{task}{OCSP}
  OpenSSL's \href{https://docs.openssl.org/master/man1/openssl-ocsp/}{\texttt{ocsp}} tool can create requests, send them, and verify the responses.
  To create the request, you need the certificate and the certificate of its issuer.
  The request does not contain the certificate itself, only a \emph{CertID}: the hash of the issuer's name, the hash of the issuer's public key, and the serial number.

  \textit{Task.}
  Export one of the certificates from your Estonian ID card.

  \begin{enumerate}
    \item Find the OCSP responder URL from the certificate.

    \item Send the OCSP request and verify the response:
    \begin{cmdblock}
openssl ocsp -issuer esteid2018.der.crt -cert cert.pem \
  -url http://aia.sk.ee/esteid2018 \
  -CAfile EE-GovCA2018.der.crt \
  -reqout req.der -respout resp.der \
  -resp_text
    \end{cmdblock}
    Use the URLs and certificates of your own certificate.

    \item Inspect the saved request and response:
    \begin{cmdblock}
openssl ocsp -reqin req.der -req_text
openssl ocsp -respin resp.der -resp_text -noverify
    \end{cmdblock}
    and verify the saved response again offline:
    \begin{cmdblock}
openssl ocsp -reqin req.der -respin resp.der \
  -CAfile EE-GovCA2018.der.crt
    \end{cmdblock}
  \end{enumerate}

  NB! Something important remains to be checked\dots.
\end{task}

\vspace*{1em}
\begin{tcolorbox}[title=Good practice]
  Chain validation and revocation checking complement each other.
  They do not replace one another!
  \begin{itemize}
    \item A CRL or OCSP response only states that an issuer has not revoked a given serial number, not that this issuer issued the certificate you are checking (e.g. a fake certificate with a copied serial).
    \item A valid chain only proves that the certificate was issued, not that it has not been revoked since.
  \end{itemize}
  First validate the chain up to a trust anchor obtained from a trusted source, then check the revocation status of every certificate in the chain.
  \tcblower
  The \texttt{CA Issuers} URL is chosen by whoever created the certificate, and is usually plain HTTP.
  Use it only for intermediates, never for roots / trust anchors.
\end{tcolorbox}
\end{document}
